\par\noindent\SourceLeaf{50}%
\hspace*{\parindent}Da $\cos\varphi^{2m-1}$ eine Reihe mit dem Anfangsglied
\[
\frac{1}{2^{2m-2}}\,{}^{2m-1}\mathfrak{B}^{(m)}\cos\varphi
\]
ist, erhält man durch Einsetzen von $90^0-\varphi$ anstelle von $\varphi$ für $\sin\varphi^{2m-1}$ das Anfangsglied
\[
\frac{1}{2^{2m-2}}\,{}^{2m-1}\mathfrak{B}^{(m)}\sin\varphi
\]
und allgemein durch dieselbe Einsetzung für $\sin\varphi^{2m-1}$ das Glied $r^{\text{tus}}$
\[
= \frac{1}{2^{2m-2}}\,{}^{2m-1}\mathfrak{B}^{(m+r)}\cos(2r+1)(90^0-\varphi).
\]
Nun gilt
\[
\begin{aligned}
\cos(2r+1)(90^0-\varphi) &= \cos\big((2r+1)90^0 - (2r+1)\varphi\big)\\
&= \cos(2r+1)90^0 . \cos(2r+1)\varphi\\&\quad + \sin(2r+1)90^0 . \sin(2r+1)\varphi.
\end{aligned}
\]
und da
\[
\sin(2r+1)90^0 = (-1)^r, \qquad \cos(2r+1)90^0 = 0,
\]
wird diese Größe zu
\[
= (-1)^r\sin(2r+1)\varphi,
\]
und das allgemeine Glied von $\sin\varphi^{2m-1}$, nämlich $r^{\text{tus}}$, zu
\[
= \frac{(-1)^r}{2^{2m-2}}\,{}^{2m-1}\mathfrak{B}^{(m+r)}\sin(2r+1)\varphi.
\]

Ersetzt man dagegen in der Reihe, die $\cos\varphi^{2m}$ darstellt, den Winkel $\varphi$ durch $90^0-\varphi$, so bleibt das Anfangsglied unverändert. Das allgemeine Glied $r^{\text{tus}}$, das
\[
\frac{1}{2^{2m-1}}\,{}^{2m}\mathfrak{B}^{(m+r)}\cos 2r\varphi,
\]
lautet, ändert sich so, dass es
\[
\begin{aligned}
\cos 2r(90^0-\varphi) &= \cos(2r\,90^0 - 2r\varphi)\\
&= \cos 2r\,90^0 . \cos 2r\varphi + \sin 2r\,90^0 . \sin 2r\varphi
\end{aligned}
\]
enthält. Nun gilt
\[
\sin 2r\,90^0 = 0, \qquad \cos 2r\,90^0 = (-1)^r,
\]
und so geht jener Ausdruck in $(-1)^r\cos 2r\varphi$ über, und für $\sin\varphi^{2m}$ wird das Glied $r^{\text{tus}}$ zu
\[
\frac{(-1)^r}{2^{2m-1}}\,{}^{2m}\mathfrak{B}^{(m+r)}\cos 2r\varphi.
\]
